\subsection{MAS-12: the character quotient metric still prevents selfadjoint similarity}
Retain the fixed character, centered quotient, weight and central removal of \fid{MAS-11}. In particular,
\[
 w_\alpha(u)=e^{2\alpha|u|}(1+u^2)^{\nu/2},\quad
 \alpha>\tfrac12,\quad\nu>1,\quad
 E_\chi=H_\alpha/V_\chi,
\]
and $R_\chi$ acts on $E_\chi^\circ=(I-P_{\chi,0})E_\chi$. All zeros below are distinct, with their multiplicities denoted by $N_\rho$. The same conclusions hold at $\alpha=1/2$ under the additional boundary nonvanishing input described in \fid{MAS-11}.

\begin{proposition}
Conditional on the exact spectral-annihilator input of \fid{MAS-11}, $R_\chi$ is not boundedly similar to a selfadjoint operator. The same is true of the operator on the even exchange summand of the labeled $\chi,\chi^{-1}$ pair. If $\chi=\chi^{-1}$ globally, the internal even operator $R_{\chi,+}$ also is not boundedly similar to a selfadjoint operator.

These assertions require neither GRH nor a simplicity hypothesis. Under GRH all retained nonzero eigenvalues are positive, but these actual quotient operators are not selfadjoint or positive, and no equivalent Hilbert norm makes them selfadjoint.
\end{proposition}
\begin{proof}
\emph{The central removal in the actual metric.}
Identify $E_\chi$ isometrically with $V_\chi^\perp\subset H_\alpha$. In the paper's conjugate-linear-first convention, the Riesz representatives of the central jets are
\[
 p_j(u)=\frac{u^j}{w_\alpha(u)},\qquad 0\le j<m,
 \quad m=\operatorname{ord}_{s=1/2}\mathcal F_\chi(s).
\]
Let $C_0$ be their span, with $C_0=0$ if $m=0$. The transpose of the Riesz spectral projection $P_0=P_{\chi,0}$ is identity on these jets and zero on all noncentral jets. Their norm-dense full span therefore gives
\[
 \operatorname{Ran}P_0^*=C_0,\qquad
 E_\chi^\circ=\ker P_0=C_0^\perp\cap E_\chi.
\]
Thus if $\Pi_0$ is the \emph{orthogonal} projection onto $C_0$ in $E_\chi$, then $Q=I-\Pi_0$ is the orthogonal projection onto $E_\chi^\circ$. It is generally different from $I-P_0$. In particular the orthogonal complement removed here is $\operatorname{Ran}P_0^*$, not $\ker P_0^*$.

For a retained critical-line zero $\rho=1/2+i\gamma$, its evaluation representative before and after the removal is
\[
 v_\gamma(u)=\frac{e^{-i\gamma u}}{w_\alpha(u)},\qquad
 q_\gamma=Qv_\gamma.
\]
The latter represents the evaluation restricted to $E_\chi^\circ$ and satisfies
\[
 R_\chi^*q_\gamma=\gamma^{-2}q_\gamma.
\]
Restrictions of finite collections of noncentral jets remain independent: a combination vanishing on $\ker P_0$ would equal its composition with $P_0$, whereas the transpose of $P_0$ kills that combination.

Set
\[
 K_\alpha(t)=\int_\R\frac{e^{itu}}{w_\alpha(u)}\,du.
\]
This is real, even and continuous, $K_\alpha(0)>0$, and $K_\alpha(t)\to0$ as $|t|\to\infty$. Every polynomial times $1/w_\alpha$ belongs to $L^1$. The Fourier decay theorem consequently gives $\|\Pi_0v_\gamma\|\to0$ as $|\gamma|\to\infty$: compute its finitely many coordinates in an orthonormal basis of $C_0$, whose elements are polynomials divided by $w_\alpha$. Direct integration gives
\begin{equation}\label{eq:MasCharacterCorrelation}
 \langle q_\gamma,q_\delta\rangle
 =K_\alpha(\gamma-\delta)
  -\langle\Pi_0v_\gamma,\Pi_0v_\delta\rangle,
 \qquad
 \|q_\gamma\|^2=K_\alpha(0)-\|\Pi_0v_\gamma\|^2.
\end{equation}
In fact $K_\alpha(t)>0$ for every real $t$, by the Gamma/Gaussian proof in \fid{MAS-09} with $e^{-|u|}$ replaced by $e^{-2\alpha|u|}$. After a nonzero central removal, the subtraction in \eqref{eq:MasCharacterCorrelation} must still be retained; it need not preserve strict positivity for every individual pair.

\emph{A growth argument forcing nearby distinct ordinates in the simple case.}
Let $n_F=[F:\mathbf Q]$. The completion fixed in \fid{MAS-11}, the identity $L_f(s,\chi)\to1$ as real $s\to+\infty$, and Stirling's formula \cite[(5.11.1)]{dlmf511} give
\begin{equation}\label{eq:MasCharacterRealGrowth}
 \log\mathcal H_\chi(x)=n_F x\log x+O_\chi(x)
 \qquad(x\to+\infty).
\end{equation}
Here $\mathcal H_\chi(x)=|f_\chi(x)|^2/|a_\chi|^2>0$ for $x>1/2$. To verify the leading coefficient, each real-place gamma factor contributes $(x/2)\log x+O_\chi(x)$ to $\log|\Lambda_\chi(1/2+x)|$, and each complex-place factor contributes $x\log x+O_\chi(x)$. Their sum is $(n_F/2)x\log x+O_\chi(x)$; taking the squared modulus doubles it. The conductor adds only a linear term, and pole clearing and central-zero removal add $O_\chi(\log x)$. The limit of the finite Euler product follows by dominated convergence in its absolutely convergent Dirichlet series for real $s\ge2$.

Suppose temporarily that every noncentral zero is simple and that each centered zero lies on the real or imaginary axis. There are only finitely many noncentral real-axis zeros, since they lie in the compact interval $[-1/2,1/2]$ and the completed entire function is nonzero. The product proved in \fid{MAS-11} yields, for large real $x$,
\[
 \log\mathcal H_\chi(x)
 =O_\chi(\log x)
  +\sum_{\gamma}\log(1+x^2/\gamma^2),
\]
where the sum ranges over its nonzero imaginary-axis zeros $i\gamma$. The finite real-axis factors contribute their logarithmic absolute values; they do not change this equality.

Let $N_{\rm im}(T)=\#\{\gamma:0<|\gamma|\le T\}$. If $N_{\rm im}(T)=O(T)$, enlarge the constant so that $N_{\rm im}(t)\le Ct$ for all $t>0$, using the zero-free neighborhood of the removed central point. Tonelli applied to
\[
 \log(1+x^2/\gamma^2)
 =\int_{|\gamma|}^{\infty}\frac{2x^2}{t(t^2+x^2)}\,dt
\]
would give
\[
 \sum_\gamma\log(1+x^2/\gamma^2)
 \le2Cx^2\int_0^\infty\frac{dt}{t^2+x^2}=\pi Cx,
\]
contradicting \eqref{eq:MasCharacterRealGrowth}. Thus $N_{\rm im}(T)$ is not $O(T)$. Since these zeros are simple, there exist distinct ordinates $\gamma_j,\delta_j$ of the same sign such that
\[
 |\gamma_j|,|\delta_j|\longrightarrow\infty,
 \qquad |\gamma_j-\delta_j|\longrightarrow0.
\]
Indeed, an eventual positive lower gap on both half-lines would give $N_{\rm im}(T)=O(T)$; discreteness permits the chosen pairs to escape each compact interval. Their squares are distinct. This step is used only after multiplicities have been disposed of below.

\emph{Exhausting the spectral and multiplicity cases.}
If some centered zero $z_\rho$ has both nonzero real and imaginary parts, then $-z_\rho^{-2}$ is a nonreal eigenvalue of $R_\chi$. Similarity preserves spectrum, so similarity to a selfadjoint operator is impossible.

Otherwise every centered zero lies on one of the two axes. If a noncentral zero has multiplicity at least two, its first two restricted jets satisfy
\[
 R_\chi'\ell_{\rho,0}=F(z_\rho)\ell_{\rho,0},\qquad
 R_\chi'\ell_{\rho,1}=F(z_\rho)\ell_{\rho,1}
                          +F'(z_\rho)\ell_{\rho,0},
 \quad F(z)=-z^{-2},\quad F'(z)=2z^{-3}\ne0.
\]
The independent jets therefore give a genuine Jordan chain at the real eigenvalue $F(z_\rho)$. Through the Riesz identification, this is a Jordan chain for the adjoint. A selfadjoint operator has no such chain: if $(A-\lambda)y=x$ and $(A-\lambda)x=0$ with real $\lambda$, then $\|x\|^2=\langle x,(A-\lambda)y\rangle=0$. Bounded similarity preserves this obstruction.

In the remaining case, all noncentral zeros are simple and lie on the axes, so the preceding growth argument gives the nearby pairs. By \eqref{eq:MasCharacterCorrelation},
\[
 \|q_{\gamma_j}\|,\|q_{\delta_j}\|\to\sqrt{K_\alpha(0)},
 \qquad \langle q_{\gamma_j},q_{\delta_j}\rangle\to K_\alpha(0).
\]
The normalized vectors $a_j,b_j$ therefore satisfy $\|a_j-b_j\|\to0$, while being eigenvectors of $R_\chi^*$ for distinct real eigenvalues.
If $R_\chi^*=SAS^{-1}$ with $A$ selfadjoint and $S$ boundedly invertible, then $S^{-1}a_j$ and $S^{-1}b_j$ are orthogonal, each of norm at least $1/\|S\|$. Their difference has norm at least $\sqrt2/\|S\|$, but also at most $\|S^{-1}\|\|a_j-b_j\|\to0$, a contradiction. Taking adjoints transfers the no-similarity conclusion from $R_\chi^*$ to $R_\chi$.

\emph{Exchange parity and internal self-dual parity.}
The even summand of the labeled pair is unitarily identified with $E_\chi^\circ$ by $x\mapsto(x,Jx)/\sqrt2$, as proved in \fid{MAS-11}. Its operator and its adjoint are therefore unitarily conjugate to $R_\chi$ and $R_\chi^*$, so the conclusion transfers with the actual metrics.

If $\chi=\chi^{-1}$ globally, reflection acts internally. Let $E_{\chi,+}$ be the original even summand, and let $C_{0,+}$ be the span of the even central vectors $u^{2j}/w_\alpha$ for $2j<m$. The orthogonal projection $Q_+$ from $E_{\chi,+}$ onto $E_{\chi,+}^\circ$ is identity minus projection onto $C_{0,+}$. For a positive critical ordinate, put
\[
 c_\gamma(u)=\frac{\cos(\gamma u)}{w_\alpha(u)},\qquad
 q_{\gamma,+}=Q_+c_\gamma.
\]
These are adjoint eigenvectors for $\gamma^{-2}$, and
\[
 \langle c_\gamma,c_\delta\rangle
 =\tfrac12\bigl(K_\alpha(\gamma-\delta)+K_\alpha(\gamma+\delta)\bigr).
\]
Their finite central projections again tend to zero as the ordinate grows. The nonreal-eigenvalue case persists on the internal even summand, because each noncentral $\{z,-z\}$ pair retains its full multiplicity there. In the multiple-zero case the independent even jets
$\ell_{z,0}+\ell_{-z,0}$ and $\ell_{z,1}-\ell_{-z,1}$ give the same nonzero Jordan entry $F'(z)$; central removal does not affect their independence.

In the remaining simple-axis case, self-duality makes the imaginary ordinates symmetric. If the positive ordinates had an eventual positive lower gap, their total two-sided count would be $O(T)$, contrary to the growth argument. Choose nearby positive pairs tending to infinity. For these pairs the displayed cosine formula and Fourier decay give
\[
 \|q_{\gamma,+}\|^2\longrightarrow K_\alpha(0)/2,\qquad
 \langle q_{\gamma,+},q_{\delta,+}\rangle\longrightarrow K_\alpha(0)/2.
\]
Normalized eigenvectors again coalesce at distinct eigenvalues, giving exactly the same bounded-similarity contradiction.

This exhausts all cases without GRH. Under GRH the exact eigenvalues from \fid{MAS-11} are positive, but selfadjointness is already excluded. A positive bounded operator is selfadjoint. Finally an equivalent Hilbert inner product is represented by a bounded positive invertible metric operator; its square root would give a bounded similarity to a selfadjoint operator if the original operator were selfadjoint in that metric. Thus an equivalent Hilbert norm cannot remove the obstruction.
\end{proof}

The ordinary character determinants of \fid{MAS-11} consequently coexist with this metric obstruction. Whenever the relevant zero-location statement is available---the stipulated theorem in the zeta/Dirichlet setting, and GRH as an additional assumption for a general Hecke character---a positive diagonal realization or the positively weighted convolution construction of \fid{DET-15} cannot be related to these weighted theta-quotient operators by a boundedly invertible intertwiner. Their determinant agreement concerns the complete algebraic spectrum and trace-class products, and does not identify the Hilbert metrics.
